\subsection{Построение функции отклонения}

В рамках лабораторной работы воспользуемся \textit{критерием Пирсона} в качестве критерия согласия.

Чтобы построить функцию отклонения \( \rho \) необходимо предварительно сгруппировать выборку. Опираясь на гипотезу \( H_1 \), определим девять квантилей для сетки \textit{(она должна быть равномерной с точки зрения вероятности)}:
\[
    \mathcal{F}(x)=
    \begin{cases}
        1-e^{-\lambda x}, &x\geq 0\\
        0, & x<0
    \end{cases}
\]

Вычислим значения квантилей при помощи скрипта на Python:

\begin{lstlisting}[caption=Расчёт квантилей экспоненциального распределения]
lmbda = 1.41

def quantile(x):
  return -math.log(1-x) / lmbda

k = 10
for i in range(1, k):
  print(i / k, f"{quantile(i / k):.4f}")
\end{lstlisting}

Итого имеем:

\begin{table}[H]
\centering
\begin{tabular}{|c|c||c|c||c|c|}
\hline
$i$ & $c_i$ & $i$ & $c_i$ & $i$ & $c_i$ \\ \hline\hline
0.1 & 0.0747 & 0.4 & 0.3623 & 0.7 & 0.8539 \\ \hline
0.2 & 0.1583 & 0.5 & 0.4916 & 0.8 & 1.1414 \\ \hline
0.3 & 0.2530 & 0.6 & 0.6499 & 0.9 & 1.6330 \\ \hline
\end{tabular}
\caption{Узлы равномерной сетки}
\label{tab:final_quantiles}
\end{table}

Определим число элементов выборки \( \nu_i \), которые попали в интервал \( \Delta_i \), при помощи скрипта:

\begin{lstlisting}[caption=Расчёт частотного распределения выборки]
X = [0.32, 0.02, 0.14, 0.52, 0.14, 1.62, 0.28, 0.22, 0.07, 0.73,
     0.83, 0.26, 1.09, 0.22, 1.43, 0.26, 0.25, 0.23, 0.05, 0.55,
     0.23, 1.57, 0.06, 1.10, 0.39, 0.49, 0.34, 1.48, 0.06, 0.52,
     0.31, 0.37, 0.15, 0.36, 0.45, 0.89, 0.35, 0.45, 0.15, 0.27,
     0.40, 0.61, 0.11, 0.06, 0.13, 0.07, 0.18, 0.03, 0.66, 2.23,
     0.04, 0.69, 0.06, 0.03, 0.96, 0.73, 0.67, 0.56, 0.02, 0.52,
     0.35, 0.09, 0.16, 0.27, 0.41, 0.21, 0.37, 0.57, 0.31, 0.39,
     0.30, 0.23, 0.56, 0.08, 0.14, 0.02, 0.58, 0.95, 0.49, 0.00,
     0.54, 1.65, 1.63, 0.23, 0.67, 0.25, 0.00, 0.23, 0.43, 0.01,
     0.10, 0.12, 2.04, 0.19, 0.07, 1.04, 0.17, 0.21, 0.55, 0.97]

grid = [0.0]
for i in range(1, k):
  grid.append(quantile(i / k))
grid.append(float('inf'))

frequencies = [0 for _ in range(k)]
for i in range(k):
  a, b = grid[i], grid[i + 1]
  for x in X:
    if a <= x < b:
      frequencies[i] += 1
print(frequencies)
\end{lstlisting}

Итого имеем:

\begin{table}[H]
\centering
\begin{tabular}{|c|c|c||c|c|c|}
\hline
$i$ & $\Delta_i$ & $\nu_i$ & $i$ & $\Delta_i$ & $\nu_i$ \\
\hline \hline
1   & $[0, 0.0747)$       & 17 & 6   & $[0.4916, 0.6499)$ & 11 \\
\hline
2   & $[0.0747, 0.1583)$  & 11 & 7   & $[0.6499, 0.8539)$ & 7  \\
\hline
3   & $[0.1583, 0.2530)$  & 15 & 8   & $[0.8539, 1.1414)$ & 7  \\
\hline
4   & $[0.2530, 0.3623)$  & 13 & 9   & $[1.1414, 1.6330)$ & 5  \\
\hline
5   & $[0.3623, 0.4916)$  & 11 & 10  & $[1.6330, +\infty)$ & 3  \\
\hline
\end{tabular}
\caption{Частотное распределение по интервалам \( \Delta_i \)}
\end{table}

На основе полученных частот посчитаем \textit{статистику Пирсона}, то есть функцию отклонения \( \rho \):
\[
    \rho(\vec{X})=\sum_{i=1}^{10}\dfrac{(\nu_i-np_i)^2}{np_i}
\]

Так как сетка равномерна с точки зрения вероятности, \( p_1=\dots=p_{10}=0.1 \). Посчитаем статистику при помощи скрипта:

\begin{lstlisting}[caption=Расчёт статистики Пирсона]
rho = 0
n = len(X)
p = 1 / k
for i in range(k):
  rho += (frequencies[i] - n * p) ** 2 / (n * p)
print(rho)
\end{lstlisting}

Итого имеем:
\[
    \rho(\vec{X})=17.8
\]

Статистика Пирсона является функцией отклонения, потому что обладает её свойствами. Это утверждает доказанная теорема:

\begin{theorem}[Пирсона]
Если верна гипотеза \( H_1 \), то при фиксированном числе групп \( k \) справедливо:
\[
    \rho(\vec{X})\xrightarrow[n\to\infty]{}H_{k-1},
\]

где \( H_{k-1} \) --- \( \chi^2 \)-распределение с \( k-1 \) степенью свободы.
\end{theorem}

Осталось определить константу \( C \). Пусть \( \eta\sim H_{k-1} \), тогда верно:
\[
    0.05=P\{\abs{\eta}\geq C\}
\]

Из таблицы распределения \( \chi^2 \) с \( 9 \) степенью свободами находим:
\[
    C=16.92
\]

Значит, критерий согласия имеет вид:
\[
    \delta(\vec{X})=
    \begin{cases}
        H_1, &\rho(\vec{X})<16.92\\
        H_2, &\rho(\vec{X})\geq 16.92
    \end{cases}
\]

Получаем, что \( \delta(\vec{X})=H_2 \), то есть выполняется \textit{альтернативная гипотеза}:
\[
    \mathcal{F}\neq\mathrm{Exp}(\lambda)
\]